\section*{Chapter \ref{ch:b3:credible-empirical-work} --- Credible Empirical Work}

\begin{solution}{exo:b3:credible-empirical-work:1}
0.647, in parentheses below 1.679. % output: reading_table
$t = 1.679/0.647 = 2.60$; the interval is $1.679 \pm 1.96 \times 0.647 = 1.679 \pm
1.268$, from 0.41 to 2.95 thousand dollars. The table should state the kind of
standard error (robust, clustered and at which level, or classical) and the
controls in each column.
\end{solution}

\begin{solution}{exo:b3:credible-empirical-work:2}
$1 - 0.95^5 = 0.226$; $1 - 0.95^{20} = 0.642$. The sixteen paths use the same data
and the same treatment variable, and the two outcomes share a component: when
one path is significant, its neighbours tend to be too, and when one is not,
neither are they. Correlated tests add fewer independent chances, so at least one
success comes less often (0.214 against 0.560). % output: forking_paths
\end{solution}

\begin{solution}{exo:b3:credible-empirical-work:3}
$\beta/s = 0.05/0.02 = 2.5$; power $= 1 - \Phi(1.96 - 2.5) = \Phi(0.54) = 0.705$.
\end{solution}

\begin{solution}{exo:b3:credible-empirical-work:4}
A symmetric inverted funnel centred on the true effect: precise studies near
it at the top, imprecise ones spread evenly on both sides at the bottom. A funnel
can also lean if small studies differ in their true effects --- for example, if
small pilot studies are run where the programme is most promising, or treat
more intensively than large roll-outs.
\end{solution}

\begin{solution}{exo:b3:credible-empirical-work:5}
Power 0.2 means $1 - \Phi(1.96 - \beta/s) = 0.2$, so $1.96 - \beta/s = 0.8416$ and
$\beta/s = 1.1184$. With $a = 0.8416$, $\lambda(a) = 0.2800/0.2 = 1.400$. The mean
significant $z$ is $1.1184 + 1.400 = 2.518$, and the ratio is $2.518/1.1184 =
2.25$, as the chart shows. % figdata: exaggeration
\end{solution}

\begin{solution}{exo:b3:credible-empirical-work:6}
The tested solution is \code{ex06_fat_pet.py}, in the chapter's \code{exercises/}
folder.
\begin{enumerate}
\item The simple mean is 0.225 and the inverse-variance mean 0.080.
% output: ex06_fat_pet
\item FAT: 2.032 (s.e.\ 0.105), $t = 19.4$; PET: 0.0056 (s.e.\ 0.0050), $t = 1.1$.
Every estimate is about twice its standard error, which is what a literature
where only significant results survive looks like. The corrected effect is
indistinguishable from zero: the literature gives no evidence of an effect.
\end{enumerate}
\end{solution}

\begin{solution}{exo:b3:credible-empirical-work:7}
With the analysis fixed in advance, only one path is reported, and its false
positive rate is the 5~\% the test promises (6.4~\% in the simulation, within
sampling error of 5~\%). % output: forking_paths
It does not change the data, the power of the study, or the filter at the
journal: a pre-registered null result can still go unpublished, which is why
registration of the study matters as much as the plan.
\end{solution}

\begin{solution}{exo:b3:credible-empirical-work:8}
Weights $1/0.01 = 100$, $1/0.0025 = 400$, $1/0.04 = 25$, sum 525. Mean $(0.20 \times
100 + 0.10 \times 400 + 0.40 \times 25)/525 = 70/525 = 0.133$; standard error
$1/\sqrt{525} = 0.044$. The second study, with the smallest standard error, has
$400/525 = 76~\%$ of the weight.
\end{solution}

\begin{solution}{exo:b3:credible-empirical-work:9}
It shows that the original estimate was probably too large: either a false
positive or, more likely in a literature with low power, a real effect
exaggerated by the filter. It does not show that the effect is zero (a third
of the original may still be real and worth knowing), that the original authors
erred, or, if the setting differs, that the original effect was not real where
it was measured.
\end{solution}

\begin{solution}{exo:b3:credible-empirical-work:10}
\begin{enumerate}
\item At $\beta/s = 1$: $a = 0.96$, $\phi(0.96) = 0.2516$, $1 - \Phi(0.96) = 0.1685$,
$\lambda = 1.493$, ratio $1 + 1.493 = 2.49$. At $\beta/s = 3$: $a = -1.04$,
$\phi(-1.04) = 0.2323$, $1 - \Phi(-1.04) = 0.8508$, $\lambda = 0.273$, ratio $1 +
0.273/3 = 1.09$.
\item As $\beta/s \to \infty$, $a \to -\infty$: $\phi(a) \to 0$ and $1 - \Phi(a) \to 1$,
so $\lambda(a) \to 0$, and $(s/\beta)\lambda(a) \to 0$; the ratio tends to 1.
\item Publication has probability $1 - \Phi(0.96) + \Phi(-2.96) = 0.1685 + 0.0015 =
0.170$. \omterm{def:b3:credible-empirical-work:exaggeration}{Sign errors} are the significant negative estimates: $0.0015/0.170 =
0.009$, about one published estimate in 110.
\end{enumerate}
\end{solution}

\begin{solution}{exo:b3:credible-empirical-work:11}
The tested solution is \code{ex11_random_effects.py}, in the chapter's
\code{exercises/} folder.
\begin{enumerate}
\item The weights are 400, 625, 100, 277.8, 400, 156.2, 1{,}111.1 and 204.1,
summing to 3{,}274.2; the weighted sum of estimates is 316.24, so $\bar\beta =
0.0966$, with standard error $1/\sqrt{3{,}274.2} = 0.0175$; $Q = 16.09$ on 7 degrees
of freedom. % output: ex11_random_effects
\item $\sum w_i^2/\sum w_i = 2{,}098{,}417/3{,}274.2 = 640.9$, so the denominator is
$3{,}274.2 - 640.9 = 2{,}633.3$ and $\hat\tau^2 = (16.09 - 7)/2{,}633.3 = 0.00345$
($\hat\tau = 0.059$). The random-effects mean is 0.1128, standard error 0.0289. The
standard error is larger because each study's variance now includes the
spread of true effects across settings, which more studies of the same size
would not remove; and the weights are more equal, so the imprecise studies,
whose estimates are larger, count for more and the mean rises.
\end{enumerate}
\end{solution}

\begin{solution}{exo:b3:credible-empirical-work:12}
An example of an answer. The training experiment of
\cref{ch:b3:difference-in-differences}: grade 2, one randomised trial in one
setting; a \omterm{def:b3:credible-empirical-work:replication}{replication} in another country would raise it. The effect of
election victory on the next election in \cref{ch:b3:regression-discontinuity}:
grade 2 for the United States House, from a strong discontinuity design; the
same design in other countries' elections, agreeing, would raise it to 1. The
effect of trade on income in \cref{ch:b3:instrumental-variables}: grade 3, since
the instrument's exclusion is disputed; a natural experiment that changed trade
alone would raise it.
\end{solution}

\begin{solution}{pb:b3:credible-empirical-work:twenty}
\textbf{1.} Nine: S02, S04, S07, S13, S14, S16, S17, S18 and S19 (for example
$21.2/7.91 = 2.68$, $3.1/1.55 = 2.00$, $3.6/1.81 = 1.99$). $9/20 = 0.45$, the share
in the summary table. % figdata: studies, studies_summary

\textbf{2.} That every study is equally informative, whatever its size, and that
the twenty are an unselected sample of the studies run.

\textbf{3.} $w_{16} = 1/1.55^2 = 0.416$; $w_{04} = 1/8.32^2 = 0.0144$; the ratio is
$8.32^2/1.55^2 = 28.8$.

\textbf{4.} The small studies report the largest effects (21.2~\%, 27.0~\%,
13.4~\%); precision weighting gives them little weight.

\textbf{5.} $3.78 \pm 1.96 \times 0.58 = 3.78 \pm 1.14$: from 2.64 to 4.92~\%.

\textbf{6.} $\beta/s = 3.78/8.32 = 0.454$; power $= 1 - \Phi(1.96 - 0.454) = 1 -
\Phi(1.506) = 0.066$. A study with one chance in fifteen of being significant
reported $z = 3.25$: a lucky draw, published because it was significant.

\textbf{7.} Estimates on the horizontal axis, precision $1/s$ on the vertical.
Without selection, the points would spread evenly on both sides of the true
effect at every height, narrowing toward the top.

\textbf{8.} S01, S02, S04, S05, S08, S09, S15 and S20. Six lie above 3.78 (21.2, 27.0,
8.8, 9.0, 9.2, 13.4) and two below (1.3 and $-0.7$): the bottom of the funnel
leans right, the mark of the filter.

\textbf{9.} Estimates rise by 1.26 for each unit of standard error; $t = 1.26/0.48
= 2.63$, so funnel asymmetry is significant at 5~\%.

\textbf{10.} The effect extrapolated to a perfectly precise study is 0.93~\%,
with $t = 0.93/1.23 = 0.76$ and an interval from $0.93 - 2.41 = -1.48$ to $0.93 +
2.41 = 3.34$~\%: corrected for selection, the twenty studies cannot tell a small
positive effect from none.

\textbf{11.} The filter acts through sampling error: a study with $s = 0$ reports
the truth and needs no luck to be significant. PET is the regression line's value
there.

\textbf{12.} $\beta/s = 0.25$, $a = 1.71$, power $= 1 - \Phi(1.71) = 0.0436$;
$\lambda = 0.0925/0.0436 = 2.12$; mean significant estimate $2 + 8 \times 2.12 =
19.0$~\%, nine and a half times the truth. S02 (21.2) and S04 (27.0) are what such
studies publish.

\textbf{13.} $2.0 \pm 1.96 \times 0.60 = 2.0 \pm 1.18$: from 0.82 to 3.18~\%.

\textbf{14.} The difference is $3.78 - 2.0 = 1.78$, with standard error
$\sqrt{0.58^2 + 0.60^2} = 0.834$: $z = 2.13$. At 5~\%, the trial rejects the
inverse-variance mean.

\textbf{15.} $2.0 - 0.93 = 1.07$, standard error $\sqrt{1.23^2 + 0.60^2} = 1.369$:
$z = 0.78$. Consistent.

\textbf{16.} $\beta/s = 2/0.6 = 3.33$; power $= \Phi(3.33 - 1.96) = \Phi(1.37) =
0.915$.

\textbf{17.} It would have been reported whatever it found, with the outcome and
the analysis fixed in advance: neither the filter nor forking paths could have
shaped it.

\textbf{18.} No. It is one estimate with an interval from 0.82 to 3.18; it agrees
with PET, but it would also have agreed with a corrected estimate of 3~\%.

\textbf{19.} Weights $1/1.23^2 = 0.661$ and $1/0.60^2 = 2.778$, sum 3.439. Combined
estimate $(0.661 \times 0.93 + 2.778 \times 2.0)/3.439 = 6.171/3.439 = 1.79$~\%,
standard error $1/\sqrt{3.439} = 0.54$.

\textbf{20.} PET rests on a model of the bias (linear in the standard error) and
extrapolates beyond the data; its standard error does not include the
uncertainty about that model. The trial alone is the cleaner number.

\textbf{21.} Before the trial: grade 2 --- several randomised trials that
disagree in size --- and the funnel shows that their agreement on a positive sign
is partly the filter's work. After it: still grade 2, now resting on one strong
pre-registered design; agreement with a second large trial elsewhere would
give grade 1.

\textbf{22.} $20{,}000 \times 0.0652 = 1{,}304$ dalens a year; $20{,}000 \times 0.020 =
400$ dalens a year.

\textbf{23.} The simple mean is $6.52/2.0 = 3.3$ times the trial's estimate: it
overstates the yearly gain by 904 dalens per trainee.

\textbf{24.} Corrected for selection and combined with the pre-registered trial,
the programme raises trainees' earnings by about 1.8~\% (standard error 0.54),
against 6.5~\% in the simple mean of the published studies; the trial alone gives
2.0~\% (0.60). Grade 2 on the \omterm{def:b3:credible-empirical-work:evidence-scale}{evidence scale}.
\end{solution}
